\section{Granularidad de la representación, Bidirectionality y caché en conversaciones de varios turnos}

La sección anterior analizó la separación de parámetros; esta sección aborda las dos estructuras restantes: que el decoder solo lee la última capa del encoder y que el encoder usa bidirectional attention. Ambas mejoraron en su momento la calidad en las tareas, pero a medida que las redes se vuelven más profundas y las conversaciones más largas, los problemas de granularidad de la información y de cálculo repetido se vuelven notorios. Aquí, por primera vez, se colocan la architecture choice y el serving cost en una misma figura.

\subsection{¿Leer solo la última encoder layer crea un cuello de botella de información?}

Esta subsección se centra en el punto de lectura de la cross-attention. En las redes profundas, distintas capas suelen representar información de distinta granularidad; si incluso la primera capa del decoder solo puede acceder a la última capa del encoder, la información formal de las capas bajas podría quedar excesivamente comprimida. En los modelos actuales de algunas decenas de capas el efecto quizá sea pequeño, pero en redes extremadamente profundas este desajuste podría amplificarse.

\lecturefigure{hyung-slide-059.jpg}{Segunda suposición: el objetivo solo necesita leer la entrada ya completamente codificada}{Hyung Won Chung official deck, p.~59}

\lecturefigure{hyung-slide-060.jpg}{La cross-attention estándar solo lee las activaciones de la última capa del encoder}{Hyung Won Chung official deck, p.~60}

\lecturefigure{hyung-slide-061.jpg}{En las redes profundas, las capas inferiores y superiores suelen codificar información de distinta granularidad}{Hyung Won Chung official deck, p.~61}

\begin{knowledgebox}{Lectura de la figura: la representation hierarchy es una heurística, no una tabla semántica fija}
En las redes de visión se suele usar «bordes en las capas bajas, objetos en las capas altas» para explicar la representación jerárquica; en los modelos de lenguaje también puede aparecer una tendencia de la forma local a la semántica abstracta. Sin embargo, el significado de cada capa cambia con el objetivo de entrenamiento, las conexiones residuales y la profundidad. Lo que la figura realmente respalda es que distintas capas pueden aportar información complementaria; por lo tanto, exponer solo la última capa es una decisión estructural que debe ponerse a prueba.
\end{knowledgebox}

\teachervoice{01:02:43--01:03:25, Hyung Won dice que en los T5 de unas 24 capas que él ha usado, que la layer 1 lea la layer 24 parece no causar problemas; pero si en el futuro hubiera diez o mil veces más profundidad, él «no se sentiría del todo cómodo». Es una pregunta típica de scale extrapolation, no una falla ya observada.}

\subsection{Qué contrafactual es la Layer-aligned connection}

La subsección anterior señaló el problema de granularidad; esta construye la alternativa mínima: que la capa $\ell$ del decoder lea la capa $\ell$ del encoder. Así, cada etapa recibe una representación de entrada de profundidad similar, lo que además se acerca más a la token interaction de la misma capa en un decoder-only. El valor del contrafactual está en que permite hacer experimentos controlados, aunque al final no se adopte.

\lecturefigure{hyung-slide-062.jpg}{Tercera suposición: dentro de la entrada se necesita una bidirectional interaction all-to-all}{Hyung Won Chung official deck, p.~62}

\lecturefigure{hyung-slide-063.jpg}{Arquitectura alternativa con layer-aligned connection y causal input}{Hyung Won Chung official deck, p.~63}

\begin{importantbox}{El producto correcto del análisis estructural es un ablation axis}
De H059--H063 se obtienen dos ejes experimentales independientes: final-layer versus layer-aligned cross-attention, y bidirectional versus causal input mask. Si se cambian ambos a la vez, la diferencia de desempeño no puede atribuirse a ninguno; la derivación en cuatro pasos del curso sirve precisamente para generar ablations que puedan escalarse por separado.
\end{importantbox}

\subsection{Beneficio en calidad y costo de sistema de la Bidirectionality}

El bidirectional encoder fue muy importante en la época de BERT para tareas de comprensión difíciles, porque cada token puede aprovechar a la vez el contexto izquierdo y el derecho. Sin embargo, en la instruction tuning a gran escala, la anecdotal experience del ponente es que la diferencia entre bidireccional y unidireccional es pequeña; al mismo tiempo, las conversaciones de varios turnos exhibieron un costo de sistema evidente: un mensaje nuevo cambia el «futuro» de la entrada anterior, y la representación bidireccional tiene que volver a codificarse.

\lecturefigure{hyung-slide-064.jpg}{A escala suficiente, el beneficio de la bidirectionality puede reducirse, pero aumenta el costo en las conversaciones de varios turnos}{Hyung Won Chung official deck, p.~64}

\teachervoice{01:03:35--01:04:23, Hyung Won marca explícitamente la conclusión como highly anecdotal: en Flan-2, la diferencia entre fine-tuning bidireccional y unidireccional es pequeña. Las notas conservan este nivel de evidencia; no debe reescribirse como «se ha demostrado que la bidirectional attention es inútil».}

\subsection{Por qué la causal attention puede reutilizar la KV cache}

Esta subsección convierte la figura en una cuenta de cómputo. En un causal model, la representación de los tokens anteriores no depende de los tokens nuevos futuros, por lo que sus key/value calculados en el turno anterior pueden guardarse en caché; al añadir un mensaje solo se calcula la parte nueva. En un bidirectional encoder, los tokens anteriores pueden leer los tokens nuevos, de modo que después de añadir algo las representaciones anteriores dejan de ser válidas y, por lo general, hay que volver a codificar toda la entrada.

\lecturefigure{hyung-slide-066.jpg}{Comparación entre la recodificación bidirectional y la cache reuse unidirectional en conversaciones de varios turnos}{Hyung Won Chung official deck, p.~66}

Sea $n$ la longitud del contexto existente y $m$ la longitud añadida. Ignorando constantes y el MLP, el costo de rehacer la full attention es aproximadamente
\[
O\bigl((n+m)^2\bigr),
\]
mientras que la causal KV cache (key-value cache, que guarda los tensores key/value de cada capa para los tokens anteriores) solo necesita calcular, para las nuevas query, la attention sobre la caché anterior y sobre los tokens añadidos, con un costo incremental aproximado de
\[
O(nm+m^2).
\]
Aquí $n$ es la longitud del contexto anterior y $m$ es el número de tokens añadidos. La caché reduce el cálculo repetido, pero aumenta el uso de memoria de la GPU; el servicio con contextos largos todavía requiere gestionar la cache placement, el batching y la eviction.

\begin{knowledgebox}{Lectura de la figura: la región azul reutilizable corresponde a la invariancia causal}
En el caso unidirectional, los tokens anteriores no pueden ver el futuro, por lo que la llegada de un mensaje nuevo no cambia los key/value anteriores; en el caso bidirectional, la representación de ``How are you?'' cambia porque después aparece ``Bad''. La ventaja de eficiencia del sistema proviene de la dirección de las dependencias en el grafo de cómputo, no del nombre decoder-only en sí.
\end{knowledgebox}

\teachervoice{La explicación en clase de 01:04:31--01:05:36 conecta la calidad del modelo con el serving: una estructura que en 2018 mejoró los benchmarks puede generar costos de recodificación en la forma de producto de varios turnos. La architecture trajectory no la determina solo la accuracy, sino también el modo de interacción.}

\subsection{Conclusión y Q\&A: el siguiente cuello de botella estructural podría ser el objective}

Al terminar el análisis de arquitectura, el ponente lleva el método de vuelta a la investigación en general. Las estructuras históricas no son reliquias irrelevantes; hacer este análisis repetidamente entrena al investigador para identificar las suposiciones implícitas del problema actual y para preguntarse si pueden sustituirse por un método más general acompañado de escala. La sesión de Q\&A plantea además que el cuello de botella más fuerte en la actualidad quizá no esté en el Transformer block, sino en el objetivo de aprendizaje.

\lecturefigure{hyung-slide-067.jpg}{Conclusión: identificar la fuerza dominante y analizar la estructura adicional desde una perspectiva de scaling}{Hyung Won Chung official deck, p.~67}

\teachervoice{El cierre de 01:06:06--01:06:42 pide a los estudiantes volver a sus propios problemas y preguntarse «qué suposiciones hay que revisar, por qué, y si pueden cambiarse por un mecanismo más general y hacer scale up». Esto es más importante que recordar que ganó el decoder-only, porque el mismo método también sirve para MoE, multimodality, memory y learning objective.}

El entrenamiento por máxima verosimilitud (maximum-likelihood estimation, MLE) toma el objetivo observado como la continuation correcta en las muestras de entrenamiento:
\[
\theta^{\star}=\arg\max_{\theta}\sum_{(x,y)\in\mathcal{D}}\log p_{\theta}(y\mid x).
\]
Aquí $\mathcal{D}$ es el conjunto de datos de pares entrada-objetivo, $x$ es la condición, $y$ es el objetivo observado y $\theta$ son los parámetros del modelo. En traducción o clasificación, una sola respuesta de referencia suele proporcionar una supervisión eficaz; en tareas abiertas como «escribe un poema», tratar un único $y$ como la única señal positiva impone una suposición estructural muy fuerte.

\begin{knowledgebox}{Qué papel desempeña aquí el RLHF}
El RLHF (reinforcement learning from human feedback) primero entrena un reward model a partir de preferencias humanas y luego usa esa recompensa para optimizar la política. Hyung Won lo considera un ejemplo de «aprender la objective function»: en lugar de solo copiar un único objetivo, se hace que el modelo aprenda la calidad relativa de varias respuestas. Al mismo tiempo, enfatiza que el RLHF tampoco es lo bastante scalable, así que esto solo demuestra una dirección; no es la solución definitiva.
\end{knowledgebox}

\teachervoice{01:12:39--01:14:34, Hyung Won dice que no cree que la architecture sea el principal bottleneck actual; le inquieta más la señal de enseñanza «solo este target es correcto» que el MLE impone a las tareas abiertas. Esta sesión de Q\&A amplía toda la clase: el inductive bias que hay que retirar también puede estar oculto en el objective.}

\begin{warningbox}{La Compute trajectory no es una promesa incondicional}
A partir de 01:15:17, el ponente califica la Moore's law de red herring y enfatiza que la variable real es la compute availability; enumera las GPU, la baja precisión y la especialización del hardware, y también reconoce que la energía y la física pueden convertirse en cuellos de botella. Un plan de investigación debe hacer análisis con varios escenarios, en lugar de extrapolar mecánicamente la tendencia exponencial pasada como una garantía para el futuro.
\end{warningbox}

\subsection{Resumen de la sección}

La jerarquía de representaciones nos recuerda revisar el final-layer bottleneck; las conversaciones de varios turnos nos recuerdan considerar la attention direction junto con el costo de la caché; la sesión de Q\&A, a su vez, extiende el análisis estructural al MLE y al RLHF. El método común sigue siendo: identificar las suposiciones que el sistema actual acepta por defecto, explicar por qué fueron útiles en el regime anterior y luego usar los cambios de escala, de datos y de interacción para comprobar si todavía vale la pena conservarlas.
